\documentclass[12pt,a4paper]{ctexart}
\usepackage[top=2.5cm,bottom=2.5cm,left=2.8cm,right=2.5cm]{geometry}
\usepackage{xcolor,hyperref,enumitem,booktabs,listings}
\hypersetup{colorlinks=true,linkcolor=black,urlcolor=blue}
\pagestyle{fancy}\fancyhf{}\fancyhead[L]{李硕·前端开发工作汇报}\fancyhead[R]{转转}\fancyfoot[C]{\thepage}
\lstset{basicstyle=\small\ttfamily,breaklines=true,frame=single,backgroundcolor=\color{gray!5}}

\begin{document}

\begin{titlepage}
\centering\vspace*{3cm}
{\Huge\bfseries 前端开发技术报告}\\[0.5cm]
{\LARGE 校园二手物品交易平台——转转}\\[1.5cm]
{\Large 李硕 · 前端开发}\\[0.3cm]
{\large 华中农业大学 · 2026年7月}
\end{titlepage}

\tableofcontents\newpage

% ============================================
\section{概述}

我在项目中负责前端页面的开发和数据准备工作。具体来说：
\begin{itemize}
  \item 首页（Hero横幅、商品网格、统计分析）
  \item 商品模块（列表搜索筛选、详情展示、发布表单）
  \item 购物车（选品、改量、结算）
  \item 订单列表（买家卖家双视角切换）
  \item 管理后台（5 个管理页面 + 侧边栏布局）
  \item ECharts 数据可视化图表
  \item 155 件商品种子数据和图片生成
\end{itemize}

以下按模块逐一说明做了什么、怎么做的、为什么这么做。

% ============================================
\section{前端基础概念速览}

在看具体代码之前，先理解几个基础概念。

\textbf{什么是组件}：Vue 中每个 \texttt{.vue} 文件就是一个组件。组件就像一个乐高积木——把页面拆分成可复用的小块。比如导航栏是一个组件，商品卡片是另一个组件，它们可以组合在一起构成完整页面。组件化的好处：修改导航栏只需要改一个文件，所有页面自动更新。

\textbf{什么是 Props 和 Emit}：父组件通过 Props（属性）向子组件传递数据，就像父母给孩子东西。子组件通过 Emit（事件）向父组件发送消息，就像孩子告诉父母"我点了一个按钮"。

\textbf{什么是 v-if 和 v-for}：\texttt{v-if} 是条件渲染——"如果条件满足才显示这段 HTML"。\texttt{v-for} 是列表渲染——"对数组中的每个元素，重复这段 HTML"。例如商品网格就是用一个 \texttt{v-for="item in products"} 循环渲染出所有商品卡片。

\textbf{什么是 CSS}：CSS（层叠样式表）控制页面的外观——颜色、大小、位置、动画。CSS Grid 和 Flexbox 是两种现代布局方式，用来决定元素在页面上的排列方式。

\section{首页是怎么工作的}

\subsection{数据加载流程}

当你打开首页（\texttt{/}），页面同时发出 4 个请求：
\begin{itemize}
  \item \texttt{GET /api/v1/index/statistics} → 获取统计数据（用户数、商品数、订单数）
  \item \texttt{GET /api/v1/category/list} → 获取分类列表（6个一级 + 20个二级）
  \item \texttt{GET /api/v1/announcement/list} → 获取公告列表
  \item \texttt{GET /api/v1/product/list?page=1\&size=20} → 获取首页商品
\end{itemize}

这四个请求是并行发出的（不是等一个完成再发下一个），这样可以更快地加载页面。每个请求都有独立的 try-catch，一个失败了不影响其他。

\subsection{商品网格的响应式布局}

商品卡片排列使用 CSS Grid 布局：

\begin{lstlisting}[language=CSS]
.product-grid {
  display: grid;
  grid-template-columns: repeat(auto-fill, minmax(240px, 1fr));
  gap: 20px;
}
\end{lstlisting}

通俗解释：让浏览器自动计算一行能放下几个卡片。每个卡片最少 240px 宽，如果屏幕够宽就多放几个，屏幕窄就少放几个。\texttt{auto-fill} 表示自动填充，\texttt{1fr} 表示每个卡片平均分配剩余空间。

\subsection{时间格式化}

商品列表中的时间是绝对时间（如"2026-07-05 14:30:00"），但展示给用户时应该用相对时间。实现逻辑：

\begin{lstlisting}[language=JavaScript]
function formatTime(t) {
  const d = new Date(t)          // 解析时间字符串
  const now = new Date()          // 当前时间
  const diff = now - d            // 时间差（毫秒）
  if (diff < 60000) return '刚刚'              // < 1分钟
  if (diff < 3600000) return Math.floor(diff/60000) + '分钟前'   // < 1小时
  if (diff < 86400000) return Math.floor(diff/3600000) + '小时前' // < 1天
  return d.toLocaleDateString('zh-CN')        // 超过1天显示日期
}
\end{lstlisting}

\subsection{公告栏滚动}

公告栏水平滚动用纯 CSS 实现，不需要 JavaScript 定时器。

核心原理：复制一份公告内容放在末尾，然后让整个内容从右到左无限移动。当第一份内容完全滚出屏幕时，它无缝衔接到末尾的副本继续滚动，视觉上形成无限循环。

\begin{lstlisting}[language=CSS]
@keyframes noticeScroll {
  0%   { transform: translateX(0); }
  100% { transform: translateX(-50%); }  /* 移动一半宽度 */
}
.notice-track {
  animation: noticeScroll 40s linear infinite;  /* 40秒一圈，匀速，无限 */
}
.notice-board:hover .notice-track {
  animation-play-state: paused;  /* 鼠标悬停暂停 */
}
\end{lstlisting}

% ============================================
\section{商品列表页的筛选系统}

\subsection{筛选参数如何在 URL 中保持}

商品列表页有 5 个筛选条件：关键词、分类、价格区间、成色、排序。如果用户筛选后刷新页面，条件不应丢失。

实现方法：把所有筛选参数写入 URL query string。例如：\texttt{/product/list?keyword=iPhone\&categoryId=7\&minPrice=1000\&sort=price\&order=asc}

\begin{lstlisting}[language=JavaScript]
// 当用户改变筛选条件时，更新 URL
function applyFilters() {
  const query = {}
  if (keyword.value) query.keyword = keyword.value
  if (categoryId.value) query.categoryId = categoryId.value
  // ...
  router.push({ query })  // 更新 URL，不刷新页面
}

// 页面初始化时，从 URL 读取筛选条件
onMounted(() => {
  if (route.query.keyword) keyword.value = route.query.keyword
  if (route.query.categoryId) categoryId.value = Number(route.query.categoryId)
  // ...
  loadProducts()
})
\end{lstlisting}

\subsection{前端如何实现分页}

分页组件（Element Plus 的 \texttt{el-pagination}）维护两个关键数据：
\begin{itemize}
  \item \texttt{page}：当前第几页
  \item \texttt{size}：每页显示多少条（默认 20）
\end{itemize}

每次翻页时，前端发送 \texttt{GET /api/v1/product/list?page=2\&size=20}，后端返回该页的数据和总条数 \texttt{total}。前端用 \texttt{total} 除以 \texttt{size} 算出总页数，显示在分页组件上。

\subsection{分类级联选择器}

商品发布时需要选择分类（如"电子产品 > 手机"）。使用 Element Plus 的 Cascader 组件，数据来自 \texttt{/category/list} 接口，返回的 JSON 已经是嵌套结构（父级含 children 数组），Cascader 可以直接使用。

\begin{lstlisting}[language=JSON]
[
  { "id":1, "name":"电子产品", "children":[
    {"id":7, "name":"手机"},
    {"id":8, "name":"电脑"}
  ]},
  { "id":2, "name":"书籍教材", "children":[...]}
]
\end{lstlisting}

用户选择"手机"后，Cascader 返回的值是 \texttt{[1, 7]}（父级ID + 子级ID），提交表单时取最后一个元素 \texttt{7} 作为 \texttt{categoryId}。

% ============================================
\section{商品详情页}

\subsection{HTML 安全过滤}

商品描述可能包含 HTML（富文本）。如果不加过滤，恶意用户可以注入 \texttt{<script>alert('攻击')</script>}，这在其他用户浏览时会被执行（XSS 攻击）。

过滤实现分三步：
\begin{enumerate}
  \item \texttt{.replace(/<script[^>]*>.*?<\/script>/gi, '')} — 移除所有 \texttt{<script>} 标签及其内容
  \item \texttt{.replace(/<style[^>]*>.*?<\/style>/gi, '')} — 移除所有 \texttt{<style>} 标签及其内容
  \item \texttt{.replace(/<[^>]*>/g, '')} — 移除所有其他 HTML 标签，只保留纯文字
  \item 通过 DOM 的 \texttt{textarea.innerHTML = ...; textarea.value} 解码 HTML 实体（如 \texttt{\&amp;} → \texttt{\&}）
\end{enumerate}

\subsection{图片画廊}

商品可以有多张图片（product\_image 表）。主图显示一张大图，下方排列缩略图。点击缩略图切换主图，通过改变 \texttt{currentImage} 响应式变量的值实现。

如果图片加载失败（网络问题或 URL 失效），\texttt{<img>} 标签的 \texttt{@error} 事件会触发，显示 Emoji 占位符（如手机分类显示 📱）。

% ============================================
\section{购物车}

\subsection{全选/取消全选}

全选状态通过计算属性动态判断：

\begin{lstlisting}[language=JavaScript]
// 当所有项的 selected 都为 true 时，全选 = true
const allSelected = computed(() =>
  list.value.length > 0 && list.value.every(i => i.selected)
)

// 切换全选：把所有项设为相同值
function toggleAll() {
  const newVal = !allSelected.value
  list.value.forEach(i => i.selected = newVal)
}
\end{lstlisting}

\subsection{防抖（Debounce）}

数量调节按钮（+/-）如果用户快速连点，会发送大量网络请求。防抖策略：用户每次点击后不立即发送请求，而是等 300ms，如果 300ms 内又点击了，就重新计时。只有最后一次点击后 300ms 内没有新点击，才真正发送请求。

\subsection{结算流程}

点击"结算"→ 弹出地址选择对话框 → 用户选择或添加地址 → 确认下单 → 为每个选中商品调用 \texttt{POST /api/v1/order} → 跳转到支付页面。

地址选择对话框从 \texttt{/api/v1/user/address} 获取用户地址列表，默认选中"默认地址"。

% ============================================
\section{订单列表的双视角切换}

这是订单列表最重要的功能改进。原始代码中 \texttt{role} 被写死为 \texttt{'buyer'}，卖家永远看不到自己的订单。

修复方案：增加一个 El-Radio-Group 切换按钮——"我买的"（role=buyer）和"我卖的"（role=seller）。切换时重新请求数据：

\begin{lstlisting}[language=JavaScript]
const role = ref('buyer')  // 默认显示"我买的"

async function loadOrders() {
  const params = { role: role.value }  // 动态切换角色
  if (activeStatus.value) params.status = activeStatus.value
  const res = await orderApi.getList(params)
  list.value = res.data || []
}
\end{lstlisting}

当切换到"我卖的"并选择"待发货"标签时，页面会显示绿色提示条，引导卖家进入订单详情发货。订单卡片上还有一个突出的"📦 去发货"按钮，点击直接跳转。

% ============================================
\section{管理后台}

\subsection{整体架构}

管理后台有独立的布局（AdminLayout），使用侧边栏导航。路由配置中设置了 \texttt{meta: \{ requiresAdmin: true \}}，路由守卫检查用户角色为非管理员时重定向到首页。

\subsection{商品审核页}

这是管理后台最复杂的页面。表格展示所有商品，筛选栏支持按状态（待审核/在售/已下架/审核驳回）过滤。操作按钮包括：
\begin{itemize}
  \item 通过：调用 \texttt{PUT /api/v1/admin/product/review/\{id\}}，状态改为"在售"
  \item 驳回：同上，状态改为"审核驳回"
  \item 下架：调用 \texttt{PUT /api/v1/admin/product/\{id\}/offline}，状态改为"已下架"
  \item 编辑：弹出对话框，可修改标题、价格、成色、状态、描述
\end{itemize}

所有操作都有确认对话框（\texttt{ElMessageBox.confirm}），防止误操作。

筛选条件变化时自动重置页码到第一页：用 \texttt{watch} 监听筛选参数变化，在回调中设置 \texttt{page = 1}。

\subsection{数据统计页——ECharts 图表}

这是项目中唯一涉及数据可视化的页面，使用 ECharts 5 库。

\textbf{图表初始化流程}：
\begin{enumerate}
  \item 在模板中放置一个 \texttt{<div ref="chartRef">} 作为图表的容器
  \item 在 \texttt{onMounted} 中获取 DOM 引用，调用 \texttt{echarts.init(chartRef.value)} 创建图表实例
  \item 调用 \texttt{chart.setOption(\{...\})} 设置图表的配置和数据
  \item 在 \texttt{onUnmounted} 中调用 \texttt{chart.dispose()} 销毁实例（释放内存）
\end{enumerate}

\textbf{饼图配置要点}：
\begin{lstlisting}[language=JavaScript]
const option = {
  series: [{
    type: 'pie',
    radius: ['40%', '70%'],   // 环形图（内径40%，外径70%）
    data: [
      { name: '手机', value: 15 },
      { name: '电脑', value: 15 },
      // ...其他分类
    ],
    label: { show: true, formatter: '{b}: {c}件 ({d}%)' }
  }]
}
chart.setOption(option)
\end{lstlisting}

\textbf{折线图配置要点}：
\begin{lstlisting}[language=JavaScript]
const option = {
  xAxis: { type: 'category', data: ['7/1','7/2',...] },  // X轴日期
  yAxis: [
    { type: 'value', name: '用户数' },    // 左Y轴
    { type: 'value', name: '订单数' }     // 右Y轴
  ],
  series: [
    { name: '新注册用户', type: 'line', data: [...] },
    { name: '新订单', type: 'line', yAxisIndex: 1, data: [...] }
  ]
}
\end{lstlisting}

\textbf{内存泄漏防护}：ECharts 实例持有 DOM 引用，组件卸载后如果不手动销毁，实例会一直占用内存。在 SPA 中频繁切换页面会积累大量未释放的图表实例。解决方案是在 \texttt{onUnmounted} 中调用 \texttt{chart.dispose()} 和 \texttt{removeEventListener}。

% ============================================
\section{商品种子数据}

\subsection{数据是怎么生成的}

155 件商品数据通过 SQL INSERT 语句批量写入。每件商品包含：
\begin{itemize}
  \item 基本信息：标题、描述、价格、原价
  \item 分类：categoryId（1-25，对应 6 个一级 + 19 个二级分类）
  \item 成色：5 个枚举值之一
  \item 卖家：user\_id 随机分配给 17 个用户
  \item 统计：viewCount、favoriteCount 随机赋值
\end{itemize}

\subsection{图片是怎么生成的}

经过 6 轮方案迭代，最终采用 Python Pillow 库自定义生成图片。脚本流程：

\begin{lstlisting}[language=Python]
# 1. 从数据库读取所有商品
for pid, title, cat_id in products:
    # 2. 根据分类选择背景色
    color, label = cat_colors.get(cat_id)

    # 3. 创建 640x480 的图片
    img = Image.new('RGB', (640, 480), rgb)

    # 4. 在图片上绘制文字
    draw = ImageDraw.Draw(img)
    draw.text((320, 240), title, fill="white", anchor="mm")

    # 5. 保存为 JPG
    img.save(f"{pid}.jpg", quality=90)
\end{lstlisting}

每张图片都包含了该商品的分类标签和商品名称，确保图文 100\% 对应。图片存储在 Docker 容器的 \texttt{/data/zhuanzhuan/uploads/products/} 目录中，通过 Nginx 代理由后端提供访问。

\textbf{为什么不用外部图库}：在项目开发过程中，我们尝试了多种外部图片来源（LoremFlickr、Picsum、Unsplash、Pexels），但都存在问题——部分图片与商品内容不符、部分 URL 失效、部分服务不稳定。最终的本地生成方案虽然图片是简单的文字+色块，但它是 100\% 可靠、100\% 匹配、0 外部依赖的。

% ============================================
\section{数据库：数据是怎么组织的}

虽然不是直接写后端，但理解数据库结构对前端开发很重要。本项目的核心数据表：

\begin{itemize}
  \item \texttt{user} 表：存储用户信息。关键字段：username（用户名）、password\_hash（BCrypt 加密的密码）、role（user/admin）、status（0禁用/1启用）。有软删除（deleted\_at）。
  \item \texttt{product} 表：存储商品信息。关键字段：user\_id（卖家）、category\_id（分类）、title、price、status（待审核/在售/已下架/已售出/审核驳回）。有软删除和 FULLTEXT 全文索引。
  \item \texttt{order} 表：存储订单。关键字段：order\_no（唯一订单号）、buyer\_id、seller\_id、product\_id、status（待付款/待发货/待收货/已完成/已取消）。
  \item \texttt{cart} 表：购物车。联合唯一键（user\_id, product\_id），同一用户同一商品只保留一条。
  \item \texttt{message} 表：站内信。from\_user\_id（发送方）、to\_user\_id（接收方）、is\_read（0未读/1已读）。
  \item \texttt{notification} 表：系统通知。user\_id（接收方）、type（order/review/system）、is\_read。
\end{itemize}

其中最重要的设计是 \texttt{order} 表的状态机——5 个状态之间有严格的转换规则：待付款→待发货→待收货→已完成，中间任何一步都可以转成"已取消"。每次状态变更都会在 \texttt{order\_log} 表中记录一条日志。

% ============================================
\section{前后端如何协作}

\subsection{数据格式约定}

前后端通过 JSON 格式通信。所有接口返回统一格式：

\begin{lstlisting}[language=JSON]
{
  "code": 200,        // 状态码：200成功，4xx客户端错误，5xx服务端错误
  "message": "success",
  "data": {...}       // 实际数据
}
\end{lstlisting}

前端 Axios 响应拦截器检查 \texttt{code}，如果是 200 就返回 \texttt{data}，否则弹出错误提示。

\subsection{接口对接流程}

以"发布商品"为例：
\begin{enumerate}
  \item 前端收集表单数据（标题、分类、价格、描述、图片ID列表）
  \item 前端通过 Axios 发送 \texttt{POST /api/v1/product}，请求头自动带 Token
  \item 后端的 Spring Security 过滤器验证 Token
  \item 验证通过后，请求到达 ProductController.publish() 方法
  \item Controller 调用 ProductService.publish()，后者开启数据库事务
  \item 事务中：插入 product 表 → 更新 product\_image 表的 product\_id → 设置封面图
  \item 返回 JSON 结果给前端
  \item 前端收到成功响应，跳转到商品详情页
\end{enumerate}

\subsection{错误处理策略}

前端对网络错误做了分类处理：
\begin{itemize}
  \item \texttt{ECONNABORTED}：请求超时（15秒），提示"检查网络连接"
  \item 无 response：网络断开，提示"网络连接失败"
  \item 401：Token 过期，自动刷新 Token
  \item 403：权限不足。如果本地有 Token 说明用户可能被删除，清除登录
  \item 404：资源不存在
  \item 429：请求太频繁
  \item 5xx：服务器错误，提示"稍后重试"
\end{itemize}

% ============================================
\section{项目中的具体问题与解决}

\textbf{问题一：ECharts 图表在窗口变大变小后不变}。原因：ECharts 不会自动响应窗口尺寸变化。解决：监听 \texttt{window resize} 事件，回调中调用 \texttt{chart.resize()}。

\textbf{问题二：管理后台筛选后跳到空页}。原因：当前在第 5 页，筛选后数据只有 2 页，但页码没变。解决：watch 筛选条件变化时重置 \texttt{page = 1}。

\textbf{问题三：购物车结算创建多个独立订单}。原因：每个选中商品分别调用 \texttt{orderApi.create}，如果中间某个失败，部分订单已创建无法回滚。目前的模拟环境可以接受，但生产环境应该由后端提供批量创建接口。

\textbf{问题四：骨架屏与实际内容高度不一致导致跳动}。解决：骨架屏卡片使用与真实卡片相同的布局结构（图片区+文字区），确保高度一致。

\textbf{问题五：商品图片加载失败破坏布局}。解决：\texttt{<img>} 标签的 \texttt{@error} 事件中替换为 Emoji 占位图。

\textbf{问题六：首页 4 个请求彼此独立，其中一个失败不应影响其他}。解决：每个请求都有独立 try-catch，失败后数据为空但页面不崩溃。

% ============================================
\section{总结}

这个项目中我写的代码主要解决了以下问题：

\begin{itemize}
  \item \textbf{首页如何展示大量商品}：CSS Grid 响应式布局 + 骨架屏加载 + 4 个并行 API 请求
  \item \textbf{筛选条件如何持久化}：URL query string 同步，刷新页面筛选不丢失
  \item \textbf{购物车的状态管理}：全选/单选用 computed 属性联动，数量修改用防抖优化
  \item \textbf{卖家如何看到并处理自己的订单}：订单列表增加"我买的/我卖的"切换
  \item \textbf{管理后台的数据可视化}：ECharts 饼图+折线图，带 resize 自适应和 dispose 防泄漏
  \item \textbf{商品数据从哪里来}：155 件种子数据 SQL + Python 自定义生成图片，100\% 可靠
\end{itemize}

% ============================================
\section{深入：CSS Grid 布局原理}

\subsection{Grid vs Flexbox}

\textbf{Flexbox}：一维布局。适合排列一行（或一列）中的元素。例如导航栏水平排列、卡片内的文字和按钮垂直排列。

\textbf{Grid}：二维布局。同时控制行和列。适合整个页面的布局、商品网格等。本项目的首页商品网格使用 Grid。

\subsection{auto-fill 和 minmax 的工作原理}

\begin{lstlisting}[language=CSS]
grid-template-columns: repeat(auto-fill, minmax(240px, 1fr));
\end{lstlisting}

这行 CSS 的背后逻辑：
\begin{enumerate}
  \item 浏览器先算出容器宽度（比如 1200px）
  \item 用 \texttt{minmax(240px, 1fr)} 来计算：每个列最少 240px，如果空间够就平均分配
  \item \texttt{auto-fill} 告诉浏览器：尽可能多地创建列，填不满就留空
  \item 1200 / 240 = 5 列
  \item 如果窗口缩小到 800px：800 / 240 = 3 列（第 4 列不够 240px，换行）
\end{enumerate}

这就是为什么同一段 CSS 在宽屏和窄屏上表现不同——浏览器自动计算列数，实现了真正的响应式布局。

\subsection{绝对单位 vs 相对单位}

\texttt{px}（像素）：绝对单位，在任何屏幕上都是一样大。\texttt{fr}（fraction）：相对单位，表示"剩余空间的 N 分之一"。\texttt{1fr 2fr} 表示前者占 1/3，后者占 2/3。

\textbf{为什么用 fr 而不是 px}：如果用固定像素，屏幕变大后卡片不会自动变宽，两侧留白。用 fr，卡片自动填满可用空间。

\texttt{rem}：相对于根元素（\texttt{<html>}）的字体大小。默认 1rem = 16px。修改根元素字体大小可以全局缩放所有使用 rem 的元素。

% ============================================
\section{深入：MySQL FULLTEXT 全文索引}

\subsection{什么是全文索引}

普通的 \texttt{LIKE '\%keyword\%'} 搜索虽然能用，但是遍历整张表逐行匹配，数据量大时非常慢（O(n) 复杂度）。全文索引通过建立倒排索引（记录每个词出现在哪些行），将搜索复杂度降到接近 O(log n)。

\textbf{倒排索引的原理}：假设有三件商品：
\begin{itemize}
  \item 商品 1："全新 iPhone 手机"
  \item 商品 2："二手 小米 手机"
  \item 商品 3："iPhone 手机壳 全新"
\end{itemize}

倒排索引：
\begin{itemize}
  \item "全新" → 商品1, 商品3
  \item "iPhone" → 商品1, 商品3
  \item "手机" → 商品1, 商品2, 商品3
  \item "小米" → 商品2
  \item "二手" → 商品2
  \item "壳" → 商品3
\end{itemize}

搜索 "iPhone 手机" 时，MySQL 查到 "iPhone" 出现在商品 1 和 3，"手机"出现在商品 1、2、3。取交集得到商品 1 和 3。这比逐行扫描快得多。

\subsection{为什么用 FULLTEXT 而不是 LIKE}

\begin{itemize}
  \item \texttt{LIKE '\%iPhone\%'}：必须扫描每一行的 title 和 description 字段，计算量大
  \item \texttt{MATCH(title,description) AGAINST('iPhone')}：利用倒排索引直接定位匹配行，速度快几个数量级
  \item 全文索引还支持自然语言模式、布尔模式（+必须包含、-排除、*通配符）等高级搜索
\end{itemize}

本项目在 product 表上建立了 \texttt{FULLTEXT KEY ft\_title\_desc (title, description)}，支持对标题和描述的同时搜索。

\subsection{全文索引的代价}

索引不是免费的。每次 INSERT 或 UPDATE 都会更新索引，占用额外的磁盘空间。所以只在搜索频繁的列上建全文索引。对于像 view\_count 这种只用于排序统计的列，普通 B-Tree 索引就够用了。

% ============================================
\section{深入：ECharts 的工作原理}

\subsection{声明式 vs 命令式}

ECharts 是声明式的：你告诉它"我要一个饼图，数据是这些，颜色用这些"，它自动完成绘制。你不需要告诉它"先画一个圆，再画一个扇形，加上标签"。这类似于 Vue 的声明式模板——你说"这个 div 里显示 count"，Vue 自动处理 DOM 更新。

\textbf{声明式的优势}：开发者关注"要什么"而不是"怎么做"；框架负责底层实现和性能优化；代码更简洁、更易维护。\textbf{命令式}（如 Canvas API）：灵活但代码量大，每个像素都要手动控制。

\subsection{ECharts 的渲染管线}

当你调用 \texttt{chart.setOption(...)} 后，ECharts 内部执行以下步骤：
\begin{enumerate}
  \item 解析配置对象，验证必要字段是否存在
  \item 数据预处理：对数据进行排序、过滤、聚合
  \item 布局计算：根据图表类型和容器大小，计算每个元素的位置和大小
  \item 坐标系统建立：确定 X 轴和 Y 轴的范围、刻度、标签
  \item 图形渲染：使用 Canvas 或 SVG 引擎绘制图形
  \item 交互绑定：绑定 hover、click 等事件处理
\end{enumerate}

其中第 3 步（布局计算）是最复杂的。饼图需要计算每个扇区的起始角度和角度范围，折线图需要把所有数据点映射到像素坐标。

\subsection{为什么折线图用双 Y 轴}

本项目的数据统计页展示两个不同量级的数据：用户注册量（通常几十到几百）和订单成交量（同样几十到几百，但趋势不同）。如果共用一个 Y 轴，两条线的波动幅度可能差异很大，视觉上不友好。

双 Y 轴（左轴和右轴各有自己的刻度）让每条线都在自己的合理范围内展示，更容易看出各自的趋势变化。

\subsection{Canvas vs SVG}

ECharts 默认使用 Canvas 渲染。Canvas 和 SVG 是两种不同的 Web 图形技术：

\textbf{Canvas}：位图模式。你告诉画笔"在坐标 (100, 200) 画一个圆"，它就真的在那个位置画像素。画完后原来的圆"消失了"（只剩像素）。适合大量图形元素、频繁重绘的场景（如数据可视化）。

\textbf{SVG}：矢量模式。每个图形是一个独立的 DOM 元素（\texttt{<circle>}、\texttt{<rect>} 等）。你可以用 CSS 修改它的样式、用 JS 绑定事件。适合图形数量不多、需要交互的场景（如图标、按钮）。

ECharts 也支持 SVG 模式（通过 \texttt{renderer: 'svg'} 配置），但 Canvas 在大数据量下性能更好。对于本项目的几十个数据点，两种模式几乎没有性能差异。

\subsection{内存泄漏的机制}

JavaScript 的内存管理依赖垃圾回收（GC）。当一个对象不再被任何变量引用时，GC 会在某个时刻回收它。

ECharts 实例持有对 DOM 元素的引用（通过 \texttt{chartRef.value}）、对 Canvas 上下文的引用、对事件监听器的引用。如果不在组件销毁时手动释放这些引用，GC 无法判断这些对象是否还需要。即使页面已经切换走了，ECharts 实例仍然"活"在内存中。

在 SPA 中，用户频繁切换页面会积累大量"僵尸"图表实例，内存占用持续增长。这就是为什么必须在 \texttt{onUnmounted} 中调用 \texttt{chart.dispose()} 和 \texttt{removeEventListener}。

% ============================================
\section{深入：MySQL 数据库设计基础}

\subsection{为什么用 InnoDB 引擎}

MySQL 支持多种存储引擎，最常用的是 InnoDB 和 MyISAM。

\textbf{InnoDB}：支持事务（ACID）、行级锁、外键约束。适合需要数据一致性的业务场景。本项目全部使用 InnoDB。

\textbf{MyISAM}：不支持事务、表级锁。读取速度快但不适合频繁写入的场景。全文索引在 MySQL 5.6 之前只有 MyISAM 支持，现在 InnoDB 也支持了。

\subsection{什么是事务}

事务是一组数据库操作，它们要么全部成功，要么全部失败（原子性）。经典的例子是银行转账：从 A 账户扣款和向 B 账户加款必须同时成功或同时失败，不能出现"扣了 A 的钱但没加到 B"。

本项目在订单创建时使用事务。Spring 的 \texttt{@Transactional} 注解让方法内的所有数据库操作在一个事务中执行。任何一步失败，整个事务回滚。

\subsection{索引的原理}

数据库索引类似于书的目录。没有索引时，查找"数据结构"这本书需要从第一页翻到最后一页（全表扫描）。有了索引，先查目录找到"数据结构"在哪一页，直接翻过去。

\textbf{B-Tree 索引}（MySQL 默认索引类型）：数据按顺序存储，查询复杂度 O(log n)。适合范围查询（>、<、BETWEEN）。

\textbf{Hash 索引}：O(1) 常数时间，但只支持等值查询（=）。MySQL 的 Memory 引擎支持 Hash 索引。

\textbf{为什么不要给所有列都建索引}：每个索引都是额外存储空间，写入（INSERT、UPDATE、DELETE）时需要同时更新所有索引。索引太多会拖慢写入速度。只在查询条件经常用到的列上建索引。

\subsection{utf8mb4 vs utf8}

MySQL 的 utf8 字符集实际只支持最多 3 字节的字符（BMP 平面），不能存储 Emoji（需要 4 字节）。utf8mb4（UTF-8 for Multilingual Plane 4-byte）是真正的 UTF-8，支持所有 Unicode 字符。本项目统一使用 utf8mb4，确保中文、Emoji 都能正确存储和显示。

% ============================================
\section{深入：前端性能优化}

\subsection{路由懒加载}

如果不做懒加载，所有页面的 JS 代码在首次访问时全部下载。对于本项目 21 个页面，这会导致首次加载缓慢。路由懒加载将每个页面分离成独立的 JS 文件（chunk），只在真正访问时才下载。Vite 构建时自动处理动态 import 的分割。

\subsection{骨架屏（Skeleton Screen）}

传统的 loading spinner（转圈圈）让用户感觉"系统在忙"。骨架屏预先画出页面的大致结构（灰色方块代替图片、灰色条代替文字），让用户感觉"内容马上就来"。研究表明骨架屏能显著降低用户的感知等待时间。

实现方式：在数据加载前用 \texttt{v-if="!loaded"} 显示骨架屏的灰色占位元素；数据加载完成后 \texttt{v-if="loaded"} 切换到真实内容。

\subsection{防抖（Debounce）和节流（Throttle）}

\textbf{防抖}：在事件触发后延迟执行，如果延迟时间内又触发了，重新计时。适用场景：搜索框输入（用户停止输入后再发请求）、购物车数量修改（停止点击后再同步）。

\textbf{节流}：无论触发多少次，固定时间间隔执行一次。适用场景：滚动事件（每 100ms 检查一次）、窗口 resize。

\subsection{为什么要把图片下载到本地}

外部 URL 图片有三个问题：加载慢（依赖第三方服务器）、不稳定（URL 可能失效）、不可控（图片内容可能变化）。下载到本地后的优势：加载快（同服务器、无跨域延迟）、稳定（不受外部影响）、可控（内容确定）。CDN 虽然能加速，但对于本项目的 155 张小图片，本地存储完全够用。

% ============================================
\section{PPT 演讲文字稿（约 1 分钟）}

以下是我在答辩 PPT 第 8 页（个人分工页）的演讲内容，基于 PPT 上展示的 6 个核心贡献点展开：

\medskip
\noindent\rule{\textwidth}{0.4pt}
\medskip

各位老师好，我是李硕，在本项目中负责前端页面开发和数据准备工作。

我主要做了五个方面的工作。

第一，首页和商品模块。首页实现了 Hero 横幅、公告栏滚动轮播、分类导航和商品网格。特别说明一下商品网格——用 CSS Grid 的 auto-fill 和 minmax 函数实现了响应式布局，屏幕宽时显示更多列，窄时自动减少。商品列表页实现了关键词搜索、分类筛选、价格区间和排序功能，所有筛选条件通过 URL 参数持久化，刷新页面不会丢失。

第二，购物车和订单列表。购物车实现了全选和单选联动、数量修改防抖同步。订单列表最重要的改进是增加了买家卖家双视角切换——原始代码只显示买到的订单，我增加了切换按钮让卖家也能看到卖出的订单，待发货的订单上还有直接跳转去发货的按钮。

第三，管理后台。我开发了 5 个管理页面——用户管理、商品审核、订单管理、数据统计和公告管理。其中数据统计页使用了 ECharts 库，用饼图展示分类占比，用折线图展示每日趋势。图表做了自适应窗口大小和组件卸载时的内存释放。

第四，数据可视化方面，我特别关注了 ECharts 的内存管理。在 Vue 的单页应用中，页面切换时如果不销毁图表实例，它们会一直占用内存。我在每个图表组件的卸载生命周期中调用了 dispose 方法来释放资源。

第五，商品数据准备。我编写了覆盖全部 20 个分类的 155 件商品种子数据。商品图片经过多次方案迭代，最终采用 Python 脚本自定义生成——每张图片都显示商品名称和分类标签，确保图文 100\% 匹配，图片存储在本地服务器，不依赖任何外部服务。

以上就是我的工作内容，谢谢各位老师。

\medskip
\noindent\rule{\textwidth}{0.4pt}
\medskip

\section{深入：RESTful API 设计}

\subsection{什么是 REST}

REST（Representational State Transfer）是一套 API 设计规范。核心思想：URL 代表资源，HTTP 方法代表操作。

本项目的 API 设计遵循 REST 风格：
\begin{itemize}
  \item \texttt{GET /api/v1/product/list} —— 获取商品列表
  \item \texttt{GET /api/v1/product/123} —— 获取商品 123 的详情
  \item \texttt{POST /api/v1/product} —— 创建新商品
  \item \texttt{PUT /api/v1/product/123} —— 更新商品 123
  \item \texttt{DELETE /api/v1/product/123} —— 删除商品 123
\end{itemize}

这种设计的好处：URL 结构一致、易于理解和文档化、前后端可以独立开发（只要约定好接口格式）。

\subsection{版本号}

本项目 API 路径中包含 \texttt{/v1/}。版本号的作用：当 API 需要不兼容的变更时，可以发布 \texttt{/v2/} 版本，旧的 \texttt{/v1/} 继续保持。这样不会破坏已有的客户端。

\subsection{统一返回格式}

所有接口返回相同的 JSON 结构：\texttt{\{code, message, data\}}。前端只需要写一个响应拦截器处理所有接口的返回，不用每个接口单独判断成功/失败。

\end{document}
